\documentclass[11pt]{article}
\usepackage[T1]{fontenc}
\usepackage[utf8]{inputenc}
\usepackage{lmodern}
\usepackage[margin=1in]{geometry}
\usepackage{amsmath,amssymb,amsthm,mathtools}
\usepackage{booktabs,array,longtable}
\usepackage{enumitem}
\usepackage{hyperref}
\newcommand{\OPHPaperReleaseID}{r1381}
\newcommand{\OPHPaperReleaseDate}{April 10, 2026}
\newcommand{\OPHPaperReleaseBanner}{%
\begin{center}
\small\textbf{Paper release:} \texttt{\OPHPaperReleaseID}\quad\textbf{Released:} \OPHPaperReleaseDate
\end{center}
}


\hypersetup{colorlinks=true,linkcolor=blue,citecolor=blue,urlcolor=blue}
\setlength{\emergencystretch}{3em}
\setlength{\tabcolsep}{5pt}
\renewcommand{\arraystretch}{1.08}

\newtheorem{proposition}{Proposition}[section]
\newtheorem{definition}[proposition]{Definition}
\newtheorem{remark}[proposition]{Remark}

\title{Toward a Particle-Spectrum Derivation from Observer-Overlap Consistency\\[0.4em]
\large A Status and Dependency Note for the OPH Particle Program}
\author{B.\ M\"uller}
\date{\OPHPaperReleaseDate}

\begin{document}
\maketitle
\OPHPaperReleaseBanner

\begin{abstract}
This note starts a dedicated public mathematical surface for the OPH particle
program. Its role is deliberately narrower than the synchronized compact paper:
it records the current dependency boundary, separates calibration outputs from
theorem-level closure claims, and names the explicit missing objects in the
charged-lepton, quark, neutrino, and hadron lanes. The present public state is
as follows. The D10 electroweak branch is a calibration-sector consistency
surface tied to the pixel-area input
\[
P \equiv a_{\mathrm{cell}}/\ell_P^2 = 1.63094,
\]
not yet an exact closure theorem. The D11 Higgs/top critical-surface branch is
a secondary quantitative continuation once D10 is fixed. Charged leptons,
quarks, and neutrinos now have explicit forward artifact boundaries and local
guard tests in the repository, but they remain continuation lanes rather than
promoted theorem-level outputs. The hadron lane remains simulation-dependent
debug and systematics scaffolding only. The purpose of the note is therefore to
state the public rule cleanly: OPH currently exposes a structured particle
derivation program with sharpened constructive objects, not a completed
first-principles closure of the full particle zoo. Accordingly, several public
mass rows still disagree with experiment, in some sectors substantially, and
those mismatches should be read as evidence of open constructive work rather
than as hidden claims of completed closure.
\end{abstract}

\section{Scope}

This note is a dependency and status surface for the OPH particle program. It
does not replace the compact submission paper as the authoritative recovered-core
statement, and it does not promote any sector merely because the code tree now
contains a forward artifact for that sector.

\begin{definition}[Particle-program claim tiers]
The public particle program uses the following four-way split.
\begin{enumerate}[leftmargin=2em]
\item \textbf{Calibration sector:} an input-dependent quantitative branch used to
fix or audit the shared electroweak family once the allowed external input
\(P\) is declared.
\item \textbf{Secondary quantitative branch:} a branch that adds no new
continuous fit once the calibration sector is fixed, but still depends on
additional transport or matching structure outside the recovered core.
\item \textbf{Continuation lane:} a forward constructive lane with explicit
objects, artifact schemas, and guards, but without theorem-level closure.
\item \textbf{Simulation-dependent lane:} a lane whose outputs currently depend
on debug, toy, or systematics-oriented numerical scaffolding and therefore
cannot be promoted as a public OPH derivation claim.
\end{enumerate}
\end{definition}

\begin{proposition}[Current public particle boundary]
On the current public repository snapshot, only the D10 electroweak surface is
presented as a calibration sector and only the D11 Higgs/top surface is
presented as a secondary quantitative branch. The charged-lepton, quark,
neutrino, and hadron lanes remain below theorem-level closure.
\end{proposition}

\paragraph{Reason.}
The compact paper and the public repository ledger now separate these lanes by
construction: the mirrored code tree carries explicit artifact boundaries and
guards for the open sectors, while the status ledger tags their rows as
continuation or simulation-dependent rather than structural or recovered-core
outputs.

\section{Dependency Ledger}

The current particle program is best read as a staged map from the OPH core and
the declared external input \(P\) into increasingly ambitious quantitative
lanes. The table below is the public boundary that README, book, and downstream
publication surfaces should mirror.

\begingroup
\small
\setlength{\tabcolsep}{3pt}
\begin{longtable}{@{}>{\raggedright\arraybackslash}p{0.14\textwidth}>{\raggedright\arraybackslash}p{0.17\textwidth}>{\raggedright\arraybackslash}p{0.22\textwidth}>{\raggedright\arraybackslash}p{0.24\textwidth}>{\raggedright\arraybackslash}p{0.17\textwidth}@{}}
\caption{Current public dependency ledger for the OPH particle program.}\\
\toprule
Lane & Current public object & Immediate inputs & Smallest stated missing object & Public tier \\
\midrule
\endhead
\bottomrule
\endfoot
D10 electroweak & running-family electroweak closure artifact & compact recovered core plus declared \(P\) and the printed running/matching conventions & one coherent \(\mathrm{EWTransportKernel}_{D10}\) with scalar-provenance equality across the electroweak correction family & calibration sector \\
D11 Higgs/top & synchronized critical-surface readout candidate & D10 plus the critical-surface synchronization branch & one common \(\mathrm{CriticalSurfaceReadoutKernel}_{D11}\) carrying the shared low-scale Higgs/top readout vector & secondary quantitative branch \\
Charged leptons & \(C_e^{\wedge} \to g_e \to Y_e\) forward chain & shared charged flavor dictionary plus charged-budget transport artifact & a theorem-level shared charged absolute-scale closure collapsing the remaining norm/gluing gap & continuation lane \\
Quarks/flavor & transport, projector, odd-response, and forward Yukawa artifacts & shared charged flavor lane plus branch-generator and overlap-edge transport objects & a theorem-level family excitation evaluator and closed odd quark response law that remove the remaining constructive ambiguity & continuation lane \\
Neutrinos & capacity anchor, family response, Majorana lift, pullback metric, splittings, and PMNS gate & D6-style capacity anchor plus shared flavor basis artifacts & an OPH-only scalar evaluator and the downstream defect-weighted family needed to promote splittings and mixing honestly & continuation lane \\
Hadrons & quenched wrapper and systematics bookkeeping & legacy lattice reference script plus explicit debug/systematics parameters & non-debug, non-quenched hadron closure with publishable systematics rather than smoke/demo artifacts & simulation-dependent lane \\
\end{longtable}
\endgroup

\section{Public Code Surface}

The repository now carries a compact public mirror of the active particle code
under \texttt{code/particles/}. Its purpose is not to reproduce the full
operational sandbox; it is to expose the current mathematical lanes together
with frozen public artifacts and tests.

The kept surfaces are:
\begin{enumerate}[leftmargin=2em]
\item the \texttt{core/} compatibility backbone used by the existing public
D10/D11 predictor surface;
\item the current lane directories
\texttt{calibration/}, \texttt{flavor/}, \texttt{leptons/},
\texttt{neutrino/}, and \texttt{hadron/};
\item the frozen artifact snapshot under \texttt{code/particles/runs/};
\item the public status ledger
\texttt{code/particles/RESULTS\_STATUS.md},
\texttt{code/particles/results\_status.json}, and
\texttt{code/particles/ledger.yaml}.
\end{enumerate}

The intentionally omitted surfaces are Oracle batch archives, browser profiles,
autonomous supervisor plumbing, large run history, and the rest of the
non-public sandbox scaffolding. Active iteration continues in the sibling
\texttt{/particles} repo; this mirror exists so the public papers and summary
surfaces can link to a compact but current code snapshot.

\section{Mathematical Closure Program}

The particle program is now organized around explicit constructive objects
rather than one broad undifferentiated ``predict the whole zoo'' claim.

\subsection{D10 exactness}

The electroweak calibration lane is no longer framed as ``perhaps more digits of
\(P\) will fix everything.'' The public D10 artifact isolates the common-family
running electroweak surface and checks whether different observables imply
different effective values of \(P\). That audit reframes the remaining burden as
one coherent transport/readout object rather than raw decimal precision.

\subsection{Charged sector}

The charged-lepton and quark lanes now meet on a shared charged-budget surface.
This matters conceptually: it prevents the public code from smuggling in one
absolute scale for leptons and another for quarks while still calling both
first-principles OPH outputs. The current public charged-sector burden is
therefore a common closure object, not two independent post-hoc fit layers.

\subsection{Neutrinos}

The neutrino lane now exposes the capacity anchor, the family-response tensor,
the Majorana lift, the pullback metric, and the forward splitting artifacts as
separate public objects. That is progress in mathematical hygiene even though it
is not yet a theorem-level neutrino closure. The public point is that the
remaining blocker is named sharply enough to fail.

\subsection{Hadrons}

The hadron branch remains the one lane where the repository is openly
simulation-dependent. The public wrapper makes that dependence explicit and
stores systematics/debug artifacts, but those artifacts are not promoted as
evidence of hadron closure.

\section{Status Surfaces and Release Discipline}

The public repository now exposes three synchronized particle-status surfaces:
\begin{enumerate}[leftmargin=2em]
\item the compact public code mirror in \texttt{code/particles/};
\item the frozen status ledger in
\texttt{code/particles/RESULTS\_STATUS.md} and
\texttt{code/particles/results\_status.json};
\item the SVG overview used by the README and public summary pages.
\end{enumerate}

These surfaces are meant to stay synchronized with the compact paper and the
main README. If a lane is still tagged \emph{continuation} or
\emph{simulation-dependent} here, downstream book and web summaries should not
silently promote it.

\paragraph{Working rule.}
The public OPH particle story is currently a sharpened constructive program, not
a completed theorem-level derivation of the full particle spectrum. When a
public mass row is still numerically wrong, the intended reading is that the
named closure object for that lane is still open, not that the current number
has already earned promotion and merely awaits cosmetic tuning.

\end{document}
